% From mitthesis package
% Version: 1.12, 2025/01/31
% Documentation: https://ctan.org/pkg/mitthesis


\chapter{Introduction}

One paragraph to tell globally what will be done in this chapter. why etc...

\newcommand*{\Jpsi}{\ifpdftex\mathord{\bm{J}/\bm{\psi}}\else\mathord{\symbfit{J/\psi}}\fi} % works with either pdftex or lualatex

\section[The messengers]{The messengers}

All different kinds of messengers. Update of what have been observed so far, and what we expect to observe in the future.


%%%%%%%%%%%%%%%%%  begin figure  %%%%%%%%%%%%%%%%%%%%%%%%%%%
%\begin{figure}[t]
% sample images are from mwe package, but should be found by latex in the tex tree w/o loading that package
%\begin{subfigure}[t]{0.495\textwidth}
%{\centering{\includegraphics[alt={sample image},width=0.99\textwidth]{example-image-c.jpg}}}%
%\subcaption{\label{fig:golden} Postremo aliquos futuros suspicor, qui me ad alias litteras vocent, genus hoc scribendi.}
%\end{subfigure}
%%%%%%%% don't leave a break here
%\begin{subfigure}[t]{0.495\textwidth}
%{\centering{\includegraphics[alt={sample image},width=0.99\textwidth]{example-image-c.jpg}}}%
%\subcaption{\label{fig:golden1} Second subfigure.}%
%\end{subfigure}%
%\caption{A figure with two subfigures.\label{fig:4}}\label{fig:golden2}
%\end{figure}
%%%%%%%%%%%%%%%%%%%  end figure  %%%%%%%%%%%%%%%%%%%%%%%%%%%%%

%% note use of \ref* here to avoid placing a nested link in the table of contents
\section[Mechanisms]{Mechanisms}

physical mechanisms that produce the messengers. This is the main part of the chapter, and it is divided into three sections, each of which is further subdivided into subsections and subsubsections.

\subsubsection{Gravitational waves}
\subsection{Electromagnetic signals}
\subsubsection{Gamma rays burst}
\subsubsection{Thermal emissions}


%\newcommand*{\boldA}{\ifpdftex\mathbf{A}\else\symbfup{A}\fi}% to get the right symbol with either pdftex or unicode-math


%% This version looks nicer, but it does not produce proper html under current tagpdf (2025/10/31)
%
%\begin{equation}\label{eqn:WT1} 
%L(\boldA) = \begin{pmatrix}
%\dfrac\varphi{(\varphi_1,\varepsilon_1)}			& 0 												 & \hdotsfor{3} 							& 0 \\[4\jot]
%\dfrac{\varphi k_{2,1}}{(\varphi_2,\varepsilon_1)}	& \dfrac\varphi{(\varphi_2,\varepsilon_2)}			 & 0 										& \hdotsfor{2} & 0 \\[4\jot]
%\dfrac{\varphi k_{3,1}}{(\varphi_3,\varepsilon_1)}	& \dfrac{\varphi k_{3,2}}{(\varphi_3,\varepsilon_2)} & \dfrac\varphi{(\varphi_3,\varepsilon_3)}	& 0 & \hdotsfor{1} & 0 \\[\jot]
%\vdots 												&  													 &  & \smash{\rotatebox{15}{$\ddots$}} &  & \vdots \\[\jot]
%\dfrac{\varphi k_{n-1, 1}}{(\varphi_{n-1},\varepsilon_1)}		& \dfrac{\varphi k_{n-1, 2}}{(\varphi_{n-1},\varepsilon_2)} & \cdots        & 
%	\dfrac{\varphi k_{n-1,n-2}}{(\varphi_{n-1},\varepsilon_{n-2})}	& \dfrac{\varphi}{(\varphi_{n-1},\varepsilon_{n-1})} 		& 0 \\[4\jot]
%\dfrac{\varphi k_{n,1}}{(\varphi_n,\varepsilon_1)}				& \dfrac{\varphi k_{n,2}}{(\varphi_n,\varepsilon_2)}		& \hdotsfor{2}	&
%	\dfrac{\varphi k_{n,n-1}}{(\varphi_n,\varepsilon_{n-1})} 		& \dfrac{\varphi}{(\varphi_n,\varepsilon_n)}
%\end{pmatrix}
%\end{equation}

%\begin{equation}
%\int_{\Gamma^h_s}\frac{\partial}{\partial n_x}\int_{\Gamma^a_\xi} \mkern-3mu b(\xi) \frac{\partial E(x|\xi)}{\partial n_\xi} \,dl_\xi \Bigg\vert_{x\to s}\,dl_s = 0 
%\end{equation}


%%%%%%%%%%%%%%% begin table %%%%%%%%%%%%%%%%%% 
%\begin{table}[t]
%\caption{The error function and complementary error function\label{tab:1}}\vspace{0.25em} % adds some additional space between caption and \toprule, if desired.
%\centering{%
%\tagpdfsetup{table/header-rows={1}}
%\begin{tabular*}{0.8\textwidth}{@{\hspace*{1.5em}}@{\extracolsep{\fill}}ccc!{\hspace*{3.em}}ccc@{\hspace*{1.5em}}}
%\toprule
%\multicolumn{1}{@{\hspace*{1.5em}}c}{$x$\rule{0pt}{8pt}} &
%$\erf{x}$ &
%\multicolumn{1}{c!{\hspace*{3.em}}}{$\erfc{x}$} &
%$x$       &
%$\erf{x}$ &
%\multicolumn{1}{c@{\hspace*{1.5em}}}{$\erfc{x}$} \\ \midrule
%0.00 & 0.00000 & 1.00000 & 1.10 & 0.88021 & 0.11980 \\
%0.05 & 0.05637 & 0.94363 & 1.20 & 0.91031 & 0.08969 \\
%0.10 & 0.11246 & 0.88754 & 1.30 & 0.93401 & 0.06599 \\
%0.15 & 0.16800 & 0.83200 & 1.40 & 0.95229 & 0.04771 \\
%0.20 & 0.22270 & 0.77730 & 1.50 & 0.96611 & 0.03389 \\
%0.30 & 0.32863 & 0.67137 & 1.60 & 0.97635 & 0.02365 \\
%0.40 & 0.42839 & 0.57161 & 1.70 & 0.98379 & 0.01621 \\
%0.50 & 0.52050 & 0.47950 & 1.80 & 0.98909 & 0.01091 \\
%0.60 & 0.60386 & 0.39614 & 1.82\makebox[0pt][l]{14} & 0.99000 & 0.01000 \\
%0.70 & 0.67780 & 0.32220 & 1.90 & 0.99279 & 0.00721 \\
%0.80 & 0.74210 & 0.25790 & 2.00 & 0.99532 & 0.00468 \\
%0.90 & 0.79691 & 0.20309 & 2.50 & 0.99959 & 0.00041 \\
%1.00 & 0.84270 & 0.15730 & 3.00 & 0.99998 & 0.00002 \\
%\bottomrule
%\end{tabular*}
%}%
%\end{table}
%%%%%%%%%%%%%%%% end table %%%%%%%%%%%%%%%%%%% 

\section{Summary}


%%	Nomenclature list is optional
%
%	This environment takes four optional arguments:
%		[1] adjust space between symbol and definition
%		[2] name (heading) of the nomenclature list
%		[3] level - can be "section" or "chapter" depending on whether you
%			have one nomenclature list for whole thesis or one for each
%			chapter. 
%		[4] style. The default matches the style of [3], but you can 
%			instead choose [frontmatter] or [backmatter] if desired.
%
%	For a single-column nomenclature list, use \begin{nomenclature}
%	For a two-column nomenclature list, use \begin{nomenclature*} AND
%		**put \usepackage{multicol}** in your preamble (in the main .tex file)
%
	\newcommand*{\boldomega}{\ifpdftex\mathord{\bm{\omega}}\else\mathord{\symbfup{\omega}}\fi} % works with pdftex or lualatex
%
\begin{nomenclature*}[2em][Nomenclature for Chapter 1][section]
\EntryHeading{Roman letters}
\entry{$\mathcal{C}$}{material curve}
\entry{$\mathbf{r}$}{material position [m]}
\entry{$\mathbf{u}$}{velocity [m s$^{-1}$]}% this more cumbersome superscripting is [arguably..] better for html translation than $^{-1}$.
\EntryHeading{Greek letters}
\entry{$\Gamma$}{circulation [m$^2$ s$^{-1}$]}
\entry{$\rho$}{mass density [kg m$^{-3}$]}
\entry{$\boldomega$}{vorticity [s$^{-1}$]} 
\end{nomenclature*}
